\chapter{Deconvolution algorithms: a complementary approach to add to the toolbox in order to counterbalance the Reproducibility Crisis in Clinical Trial Research}

In all the sections described before, we applied the same biological model to the patients of a given cohort. However, we observe strong variability between the samples, even within the same individual. The variability occurs at three levels: at the environmental level (disease state, tissue location, . . . ), the genotypical level (presence of individual mutations inducing varying transcriptomic activity) and even at the cell population level. 

In the next sections, we will first detail canonical clustering methods, namely \Glspl{gmm} that are still relevant in the exploratory stage to unravel homogeneous groups in a seemingly highly variable disease condition, both at the phenotypical and transcriptomic level. Then, we review some of the most common deconvolution methods designed to measure and for some of them, address the bias induced by the strong variability of the cell composition.

\paragraph{the main sources of variability of the transcriptomic expression}

\section{Mixture models to model highly heterogeneous diseases} 

\subsection{A review of Gaussian mixture models in R}

\subsection{A practical application on the complex immuno-inflammatory Sjögren's disease}


\section{A state-of-the-art of the deconvolution approaches used to infer cell ratios, purified cellular profiles or both}

\subsection{Physical approaches used to analyse the cytometry environment}

\subsection{Reference-free approaches}

\subsection{Partial deconvolution methods}

\subsubsection{Marker-based approaches}

\subsubsection{Partial based approaches with the complete signature matrix}




